%% kapa: exam=1 round=20 year=2009 period=2 form=A subject=accounting subject_ko=회계학 seq=21 number=21
%% source: Q-Net 감정평가사 기출문제 (cst00309) artlSeq=5024178 file=extracted/1008/5024178/2009년 제20회 감정평가사 1차시험 기출문제/2교시.hwp
%% conversion: kordoc 4.24.51 -> tools/gen1.py (auto); review=proofread-sonnet; flags=none
\begin{question}{21}
(주)한국은 20×8년 7월 1일에 취득원가 ￦246,000의 기계장치를 취득하고 연수합계법으로 감가상각을 해 오고 있다. 이 기계장치의 추정내용연수는 4년이며 추정잔존가치는 ￦6,000이다. 만약 이 기계장치를 정률법에 의해서 상각하였다면 연수합계법으로 상각한 경우와 비교하여 20×9년 말의 감가상각누계액은 얼마나 차이가 있는가? (단, 정률법 감가상각에서 상각률은 50\%라고 가정하며, 감가상각은 월할계산 한다.)
\begin{choices}
  \item ￦8,250
  \item ￦15,750
  \item ￦22,000
  \item ￦21,750
  \item ￦28,500
\end{choices}
\end{question}
\answer{4}
